\section{Waymo dataset}

\subsection{Dataset Information}
The Waymo Dataset~\cite{waymo} is a high-quality multimodal dataset for autonomous driving perception. Significant for its scale and diversity, the dataset comprises 1,150 scenes, each spanning 20 seconds, collected across a wide range of urban and suburban geographies including San Francisco, Phoenix, and Mountain View. The sensor setup captures a full 360-degree view using a complex suite of 5 LiDARs (one mid-range, four short-range) and 5 high-resolution cameras. The dataset includes diverse driving conditions, covering various weather patterns (rain, fog, direct sunlight) and times of day (dawn, day, dusk, night).

The dataset contains approximately 12 million annotated 3D bounding boxes across 3 main object classes: \textit{Vehicle, Pedestrian,} and \textit{Cyclist}. Unlike the relational database structure of nuScenes, the Waymo Dataset is distributed in the TFRecord format using Protocol Buffers. The data is split into three main sets:

\begin{itemize}
    \item \texttt{training}: 798 scenes (approx. 158k frames) with annotations.
    \item \texttt{validation}: 202 scenes (approx. 40k frames) with annotations.
    \item \texttt{testing}: 150 scenes (approx. 30k frames) without published annotations.
\end{itemize}



The label annotations are encapsulated within \texttt{Frame} protocol buffers. Each 3D object label typically contains the following fields:

\begin{itemize}
    \item \texttt{type}: The object class (e.g., \texttt{TYPE\_VEHICLE}).
    \item \texttt{box\_3d}: Contains center $(x, y, z)$, dimensions $(l, w, h)$, and heading $\theta$.
    \item \texttt{metadata}: Tracking difficulty and speed estimates.
    \item \texttt{detection\_difficulty\_level}: Human-annotated difficulty (LEVEL\_1 or LEVEL\_2).
    \item \texttt{num\_lidar\_points\_in\_box}: The number of LiDAR returns inside the 3D box.
    \item \texttt{tracking\_id}: Unique identifier for tracking objects across frames.
\end{itemize}

\subsection{Evaluation Metrics}
The Waymo object detection benchmark evaluates performance using mean Average Precision (mAP) and a heading-weighted metric called mean Average Precision weighted by Heading (mAPH). Unlike nuScenes which uses center-distance matching, Waymo utilizes 3D Intersection over Union (IoU) thresholds for matching predictions to ground truth.

A prediction is matched if its IoU with a ground truth box exceeds the class-specific threshold (0.7 for vehicles, 0.5 for pedestrians and cyclists). The primary metrics are defined as follows:

\begin{itemize}
    \item \textbf{mAP}: Standard mean Average Precision based on 3D IoU.
    \item \textbf{mAPH}: Incorporates heading accuracy; a True Positive is weighted by $\frac{1 + \cos(\alpha)}{2}$, where $\alpha$ is the heading error.
\end{itemize}

To provide a granular analysis of model performance, metrics are broken down by difficulty levels based on the number of LiDAR points visible on the object:

\begin{itemize}
    \item \textbf{LEVEL\_1}: Objects with $>5$ LiDAR points. Considered "easy" to detect.
    \item \textbf{LEVEL\_2}: Objects with $\ge 1$ LiDAR point (includes LEVEL\_1). Represents the full difficulty range including occluded or distant objects.
\end{itemize}

The final mAPH score is calculated as:
\begin{equation}
    \text{mAPH} = \int_{0}^{1} p(r) \cdot h(r) \, dr
\end{equation}
\noindent where $p(r)$ is the precision at recall $r$, and $h(r)$ is the heading accuracy weight for the matched examples.